\documentclass[10pt,a4paper]{article}
\usepackage[top=1.6cm,bottom=1.4cm,left=1.2cm,right=1.2cm,headsep=8pt,footskip=16pt]{geometry}
\usepackage{enumitem}
\usepackage{amssymb}
\usepackage{array}
\usepackage[table]{xcolor}
\usepackage{fancyhdr}
\usepackage{multicol}
\usepackage{polyglossia}
\setdefaultlanguage[calendar=gregorian, numerals=maghrib]{arabic}
\setotherlanguage{english}
\IfFontExistsTF{Amiri}{%
  \newfontfamily\arabicfont[Script=Arabic, Scale=1.0]{Amiri}
}{%
  \newfontfamily\arabicfont[Script=Arabic, Scale=1.0, Path=../fonts/, Extension=.ttf,
    UprightFont=*-Regular, BoldFont=*-Bold, ItalicFont=*-Italic,
    BoldItalicFont=*-BoldItalic]{Amiri}
}
\IfFontExistsTF{AlBattar}{%
  \newfontfamily\khatt[Script=Arabic, Scale=1.15]{AlBattar}
}{%
  \newfontfamily\khatt[Script=Arabic, Scale=1.15, Path=../fonts/,
    BoldFont=ae_AlBattar.ttf]{ae_AlBattar.ttf}
}
\renewcommand{\headrulewidth}{0.8pt}
\renewcommand{\footrulewidth}{0.8pt}
\fancyhf{}
\rhead{\khatt المدرسة العليا للأساتذة ببوزريعة}
\lhead{\khatt ملحقة بومرداس}
\lfoot{\khatt الأستاذ: د. بن عبد الرزاق يوسف}
\cfoot{\thepage/4}
\pagestyle{fancy}
\parindent 0pt
\setlength{\columnsep}{0.6cm}
\setlength{\columnseprule}{0.4pt}
\setlist[enumerate,1]{label=\textbf{\arabic*.}, leftmargin=*, itemsep=5pt, topsep=2pt}
\begin{document}

\begin{center}
{\LARGE\khatt\textbf{الاختبار في الذكاء الاصطناعي}}\\[3pt]
{\normalsize أسئلة متعددة الاختيارات ــ الفصول من 1 إلى 6 ــ 20 سؤالًا}
\end{center}
\vspace{-4pt}\rule{\textwidth}{0.6pt}\\
{\small\textbf{التعليمة:} لكل سؤال إجابة صحيحة واحدة فقط من بين خمسة اختيارات؛ اقرأ الاختيارات كلها بعناية قبل الإجابة.
اكتب \textbf{حرف} الإجابة الصحيحة في جدول الإجابات أسفل الصفحة الأخيرة.
لكل سؤال نقطة واحدة (المجموع 20 نقطة).}
\vspace{-2pt}
\begin{multicols}{2}

\noindent\begin{minipage}{\linewidth}
\hangindent=0.7cm\hangafter=1\makebox[0.7cm][r]{\textbf{1.}}\textbf{في ورشة تكوينية، طُلب من المشاركين تعريف الذكاء الاصطناعي التوليدي تعريفًا يميّزه عن غيره من التطبيقات الذكية. أي تعريف هو الأدق؟}
\begin{itemize}[leftmargin=0.95cm, labelsep=4pt, itemsep=1pt, topsep=3pt, parsep=0pt]
  \item[\textbf{أ)}] نوع من الذكاء الاصطناعي ينتج محتوى جديدًا، كالنصوص والصور والأصوات، انطلاقًا من طلب يكتبه المستخدم، ولا يكتفي بإجابة جاهزة مخزنة
  \item[\textbf{ب)}] جهاز إلكتروني حديث يُركَّب في الحواسيب لتوليد الطاقة وتسريع المعالجة، وسُمّي توليديًا لأنه يولّد الكهرباء اللازمة لتشغيل البرامج الذكية
  \item[\textbf{ج)}] برنامج لتصحيح الأخطاء الإملائية والنحوية فقط، يطبّق قواعد ثابتة كتبها اللسانيون يدويًا، ولا يستطيع إنتاج الصور ولا الأصوات ولا الجداول
  \item[\textbf{د)}] قاعدة بيانات ضخمة تحفظ كل الإجابات الممكنة لكل الأسئلة، فإذا كتب المستخدم طلبه بحثت عن الإجابة المطابقة له حرفيًا وعرضتها كما هي
  \item[\textbf{هـ)}] نوع من محركات البحث يجمع الصفحات الأكثر زيارة حول موضوع الطلب، ثم يعرضها مرتبة كما هي، دون أن ينتج أي محتوى لم يكن موجودًا قبل ذلك
\end{itemize}
\end{minipage}\par\vspace{8pt plus 6pt}
\noindent\begin{minipage}{\linewidth}
\hangindent=0.7cm\hangafter=1\makebox[0.7cm][r]{\textbf{2.}}\textbf{شرح أستاذ لطلبته مراحل تعلّم الآلة مستعملًا تشبيه التلميذ الذي يتعلم الإعراب. أي وصف يعرض المراحل الخمس بترتيبها الصحيح؟}
\begin{itemize}[leftmargin=0.95cm, labelsep=4pt, itemsep=1pt, topsep=3pt, parsep=0pt]
  \item[\textbf{أ)}] ننتج النموذج أولًا، ثم نجمع البيانات لتدريبه، ثم نعرض النتيجة على المستخدم، وأخيرًا نختبره بالأمثلة نفسها التي تدرّب عليها
  \item[\textbf{ب)}] ندرّب النظام أولًا دون بيانات، ثم نجمع البيانات لاختباره، ثم ينتج النموذج، ثم تظهر النتيجة، وأخيرًا نعيد جمع البيانات من جديد
  \item[\textbf{ج)}] نختبر النظام أولًا لنعرف مستواه، ثم نجمع البيانات، ثم ندرّبه عليها، ثم نعرض النتيجة، وأخيرًا ينتج عن كل ذلك النموذج النهائي
  \item[\textbf{د)}] نجمع البيانات، ثم يتعلم النظام منها في مرحلة التدريب فينتج النموذج، ثم نختبره بأمثلة جديدة لم يرها، وأخيرًا يقدّم النتيجة
  \item[\textbf{هـ)}] نجمع البيانات، ثم نعرض النتيجة مباشرة على المستخدم، ثم ندرّب النظام على ملاحظاته، ثم ننتج النموذج، ثم نختبره بالبيانات الأولى
\end{itemize}
\end{minipage}\par\vspace{8pt plus 6pt}
\noindent\begin{minipage}{\linewidth}
\hangindent=0.7cm\hangafter=1\makebox[0.7cm][r]{\textbf{3.}}\textbf{كتب أستاذ الأمر: «تصرّف كمفتش تربوي. اقترح خمسة أنشطة لدرس المفعول به لتلاميذ السنة الأولى متوسط، في جدول». ما العنصر الذي تمثله عبارة «تصرّف كمفتش تربوي»، وعن أي سؤال يجيب؟}
\begin{itemize}[leftmargin=0.95cm, labelsep=4pt, itemsep=1pt, topsep=3pt, parsep=0pt]
  \item[\textbf{أ)}] يمثل عنصر التنسيق، ويجيب عن السؤال: كيف نريد شكل الإجابة؟ هل في جدول أم في قائمة أم في فقرات متتابعة؟
  \item[\textbf{ب)}] يمثل عنصر المهمة، ويجيب عن السؤال: ماذا نريد أن تفعل الأداة بالضبط؟ هل تقترح أنشطة أم تشرح القاعدة؟
  \item[\textbf{ج)}] يمثل عنصر الدور، ويجيب عن السؤال: من تكون الأداة؟ وأي خبرة تعتمد عليها عند صياغة الإجابة المطلوبة منها؟
  \item[\textbf{د)}] يمثل عنصر القيود، ويجيب عن السؤال: ما الشروط التي يجب أن تحترمها الأداة في عدد الأنشطة وطولها ولغتها؟
  \item[\textbf{هـ)}] يمثل عنصر السياق، ويجيب عن السؤال: ما المعلومات التي تساعد الأداة على فهم وضعية التلاميذ ومستواهم الدراسي؟
\end{itemize}
\end{minipage}\par\vspace{8pt plus 6pt}
\noindent\begin{minipage}{\linewidth}
\hangindent=0.7cm\hangafter=1\makebox[0.7cm][r]{\textbf{4.}}\textbf{يقوم الاستعمال الأخلاقي للذكاء الاصطناعي في الدرس على ثلاثة مبادئ أساسية. أي اقتراح يعرض هذه المبادئ مع معانيها الصحيحة؟}
\begin{itemize}[leftmargin=0.95cm, labelsep=4pt, itemsep=1pt, topsep=3pt, parsep=0pt]
  \item[\textbf{أ)}] المسؤولية: الأداة مسؤولة عن أخطائها؛ الشفافية: أخفي طريقة الاستعمال؛ الإفصاح: أنشر إجابة الأداة كما هي دون أي تعديل
  \item[\textbf{ب)}] المسؤولية: الأستاذ مسؤول عن أخطاء الطلبة؛ الشفافية: أستعمل أداة مجانية؛ الإفصاح: أترجم إجابة الأداة قبل نشرها
  \item[\textbf{ج)}] المسؤولية: أنا مسؤول عمّا أقدّمه؛ الشفافية: أوضّح كيف استعملت الأداة؛ الإفصاح: أصرّح باستعمالها عندما يكون ذلك مطلوبًا
  \item[\textbf{د)}] المسؤولية: لا أحد مسؤول عن أخطاء الأداة؛ الشفافية: أختار أداة مشهورة؛ الإفصاح: أحذف المصادر حتى لا يطول البحث
  \item[\textbf{هـ)}] المسؤولية: الشركة مسؤولة عن كل استعمال؛ الشفافية: أكتب بخط واضح؛ الإفصاح: أذكر اسمي فقط في أول صفحة من البحث المقدَّم
\end{itemize}
\end{minipage}\par\vspace{8pt plus 6pt}
\noindent\begin{minipage}{\linewidth}
\hangindent=0.7cm\hangafter=1\makebox[0.7cm][r]{\textbf{5.}}\textbf{قارنت أستاذة إجابتي أداتين مختلفتين عن الطلب نفسه «اشرح النعت لتلاميذ السنة الأولى متوسط مع ثلاثة أمثلة»، فوجدتهما مختلفتين في الطول والتنظيم، وفي إحداهما مثال خاطئ. ما التصرف الأنسب؟}
\begin{itemize}[leftmargin=0.95cm, labelsep=4pt, itemsep=1pt, topsep=3pt, parsep=0pt]
  \item[\textbf{أ)}] ترفض الأداتين معًا، لأن اختلاف إجابتيهما يدل على أن الذكاء الاصطناعي لا يصلح إطلاقًا لإعداد دروس اللغة العربية
  \item[\textbf{ب)}] تختار الإجابة الأطول مباشرة، لأن الطول دليل على الدقة، وتستعملها كما هي مع تلاميذها دون الحاجة إلى مراجعة الأمثلة
  \item[\textbf{ج)}] تعتمد الأمثلة المشتركة بين الأداتين فقط، لأن اتفاق أداتين على معلومة يكفي دليلًا على صحتها دون الرجوع إلى أي مصدر
  \item[\textbf{د)}] تعيد الطلب نفسه مرات عديدة حتى تتطابق الإجابتان حرفيًا، لأن تطابق الإجابات هو الدليل الوحيد على صحتها وموثوقيتها
  \item[\textbf{هـ)}] تستفيد من المقارنة لاختيار الأوضح، ثم تتحقق من كل مثال في كتاب نحو معتمد، وتصحّح الخطأ قبل استعمال أي شيء مع تلاميذها
\end{itemize}
\end{minipage}\par\vspace{8pt plus 6pt}
\noindent\begin{minipage}{\linewidth}
\hangindent=0.7cm\hangafter=1\makebox[0.7cm][r]{\textbf{6.}}\textbf{درّب فريق نموذجًا لتصنيف أخطاء التلاميذ الإملائية على 900 مثال من السنة الرابعة متوسط و100 مثال فقط من السنة الأولى. ما النتيجة المتوقعة، ولماذا؟}
\begin{itemize}[leftmargin=0.95cm, labelsep=4pt, itemsep=1pt, topsep=3pt, parsep=0pt]
  \item[\textbf{أ)}] يخطئ النموذج كثيرًا مع أخطاء تلاميذ السنة الأولى، لأن البيانات غير متوازنة ولا تمثّل هذه الفئة، وهذا شكل من الانحياز
  \item[\textbf{ب)}] يرفض النموذج تصنيف أي خطأ من السنة الأولى، لأن النماذج لا تتعامل إلا مع الفئة التي تمثل أكثر من نصف بيانات التدريب
  \item[\textbf{ج)}] يعمل النموذج بالدقة نفسها مع الفئتين، لأن العدد الكلي للأمثلة كبير، وهذا وحده كافٍ لضمان نتائج عادلة ودقيقة للجميع
  \item[\textbf{د)}] يتعرف النموذج على أخطاء السنة الأولى أفضل من غيرها، لأن الأمثلة القليلة تجعله يركّز عليها أكثر أثناء مرحلة التدريب
  \item[\textbf{هـ)}] لا يمكن توقع أي شيء قبل مرحلة النتيجة، لأن توازن البيانات لا علاقة له بجودة التعلم، فالمهم هو سرعة الحاسوب المستعمل
\end{itemize}
\end{minipage}\par\vspace{8pt plus 6pt}
\noindent\begin{minipage}{\linewidth}
\hangindent=0.7cm\hangafter=1\makebox[0.7cm][r]{\textbf{7.}}\textbf{كتب طالب الأمر: «اشرح التشبيه»، ثم أراد تحسينه باتباع البناء التدريجي للأمر الذي عرضه الدرس. أي صيغة نهائية تمثل أفضل تحسين؟}
\begin{itemize}[leftmargin=0.95cm, labelsep=4pt, itemsep=1pt, topsep=3pt, parsep=0pt]
  \item[\textbf{أ)}] «اشرح التشبيه شرحًا جيدًا جدًا ومفصلًا وواضحًا ودقيقًا وكاملًا، ولا تنسَ أي شيء، وأعطني أفضل إجابة ممكنة»
  \item[\textbf{ب)}] «التشبيه، التشبيه، التشبيه. أريد شرحًا للتشبيه الآن، وأعد كتابة الإجابة إذا لم تعجبني حتى أحصل على ما أريده»
  \item[\textbf{ج)}] «أنا طالب في السنة الأولى. اشرح التشبيه في 150 كلمة بلغة بسيطة، مع مثالين، ونظّم الإجابة: تعريف، أمثلة، خلاصة»
  \item[\textbf{د)}] «تصرّف كأستاذ. اشرح شيئًا عن البلاغة بطريقة جميلة، واكتب ما تراه مناسبًا من أمثلة أو جداول أو أي شكل تريده»
  \item[\textbf{هـ)}] «اشرح التشبيه، واكتب قصيدة فيها تشبيهات، وترجمها إلى الإنجليزية، ولخّص درس الاستعارة، واقترح امتحانًا في البلاغة»
\end{itemize}
\end{minipage}\par\vspace{8pt plus 6pt}
\noindent\begin{minipage}{\linewidth}
\hangindent=0.7cm\hangafter=1\makebox[0.7cm][r]{\textbf{8.}}\textbf{أنشأ طالب مسودة بحث بأداة ذكاء اصطناعي، ثم راجعها وتحقق من معلوماتها وأعاد صياغتها، وكتب جملة إفصاح في المقدمة. كيف يُصنَّف هذا الاستعمال حسب مقياس الدرس؟}
\begin{itemize}[leftmargin=0.95cm, labelsep=4pt, itemsep=1pt, topsep=3pt, parsep=0pt]
  \item[\textbf{أ)}] انتحال مرفوض: كل استعمال للأداة في كتابة بحث يُعدّ انتحالًا، حتى لو راجع الطالب المسودة وتحقق منها وصرّح بذلك صراحة
  \item[\textbf{ب)}] غير مسؤول: لأن كل اعتماد على مسودة من الأداة يعني بالضرورة أن الطالب لم يتحقق من معلوماتها، مهما قال إنه راجعها
  \item[\textbf{ج)}] ممنوع دائمًا: لأن المؤسسات الجامعية كلها تمنع أي استعمال للذكاء الاصطناعي في البحوث، بغض النظر عن طريقة الاستعمال
  \item[\textbf{د)}] مقبول بشروط: إنشاء مسودة ثم مراجعتها والتحقق منها والإفصاح عنها استعمال مقبول ما دام الطالب يتحمل مسؤولية النص النهائي
  \item[\textbf{هـ)}] مقبول دون أي شرط: لأن الإفصاح وحده يكفي، فلا حاجة إلى مراجعة المسودة ولا إلى التحقق من معلوماتها ولا إلى إعادة صياغتها
\end{itemize}
\end{minipage}\par\vspace{8pt plus 6pt}
\noindent\begin{minipage}{\linewidth}
\hangindent=0.7cm\hangafter=1\makebox[0.7cm][r]{\textbf{9.}}\textbf{تريد طالبة إعداد بحث حول موضوع «الذكاء الاصطناعي والتعليم». نصحها الأستاذ بالبدء بسؤال بحث، ثم باستخراج الكلمات المفتاحية. أي اقتراح يطبّق النصيحة تطبيقًا صحيحًا؟}
\begin{itemize}[leftmargin=0.95cm, labelsep=4pt, itemsep=1pt, topsep=3pt, parsep=0pt]
  \item[\textbf{أ)}] السؤال: «هل الذكاء الاصطناعي مفيد أم لا؟»؛ الكلمات: هل، مفيد، أم، لا، لأنها تمثل السؤال كما هو دون حذف أي كلمة منه
  \item[\textbf{ب)}] السؤال: «كيف يساعد الذكاء الاصطناعي التلميذ على تحسين التعبير الكتابي؟»؛ الكلمات: الذكاء الاصطناعي، التعبير الكتابي، \textenglish{Writing skills}
  \item[\textbf{ج)}] السؤال: «ما كل شيء عن الذكاء الاصطناعي؟»؛ الكلمات: كل، شيء، عن، ما، لأن الكلمات القصيرة تعطي نتائج أكثر في محرك البحث
  \item[\textbf{د)}] السؤال: لا حاجة إليه، فالموضوع يكفي؛ الكلمات: نكتب الموضوع كاملًا في محرك البحث، ونأخذ النتيجة الأولى لأنها الأفضل
  \item[\textbf{هـ)}] السؤال: «الذكاء الاصطناعي والتعليم في كل دول العالم»؛ الكلمات: كيف، ما، هل، لماذا، متى، أين، وكل الكلمات العامة الأخرى
\end{itemize}
\end{minipage}\par\vspace{8pt plus 6pt}
\noindent\begin{minipage}{\linewidth}
\hangindent=0.7cm\hangafter=1\makebox[0.7cm][r]{\textbf{10.}}\textbf{رفع أستاذ صورة مخطط بياني لنتائج قسمه إلى أداة متعددة الأنماط وطلب منها شرحه، فقرأت الأداة رقمًا خطأ وخلطت بين عمودين. ماذا يجب أن يفعل وفق خطوات الدرس؟}
\begin{itemize}[leftmargin=0.95cm, labelsep=4pt, itemsep=1pt, topsep=3pt, parsep=0pt]
  \item[\textbf{أ)}] يعتمد الشرح كما هو، لأن الأداة متعددة الأنماط تفهم الصور أفضل من الإنسان، والخطأ في الأرقام لا يغيّر المعنى العام
  \item[\textbf{ب)}] يتوقف نهائيًا عن استعمال الأداة، لأن خطأً واحدًا في قراءة الأرقام يدل على أنها لا تصلح لأي استعمال تعليمي مستقبلًا
  \item[\textbf{ج)}] يطلب من الأداة تحويل المخطط إلى تسجيل صوتي، لأن الأرقام تُقرأ بدقة أكبر عند تحويلها من نمط الصورة إلى نمط الصوت
  \item[\textbf{د)}] يحذف الأعمدة التي أخطأت فيها الأداة من المخطط الأصلي، ثم يقدّم الشرح لتلاميذه، لأن الأعمدة المتبقية صحيحة بالتأكيد
  \item[\textbf{هـ)}] يقارن الشرح بالصورة الأصلية رقمًا رقمًا، ويصحّح الأخطاء، ويمكنه إعادة رفع صورة أوضح مع أمر دقيق قبل اعتماد أي شرح
\end{itemize}
\end{minipage}\par\vspace{8pt plus 6pt}
\noindent\begin{minipage}{\linewidth}
\hangindent=0.7cm\hangafter=1\makebox[0.7cm][r]{\textbf{11.}}\textbf{قال أستاذ: «طلبتُ من الأداة نسبة بيت شعري إلى قائله، فأجابت بثقة كبيرة وبأسلوب جميل، إذن الإجابة صحيحة». أي تحليل يبيّن خطأ هذا الاستدلال بدقة؟}
\begin{itemize}[leftmargin=0.95cm, labelsep=4pt, itemsep=1pt, topsep=3pt, parsep=0pt]
  \item[\textbf{أ)}] الخطأ أنه لم يكتب الطلب بالإنجليزية، فالأدوات لا تهلوس إلا مع الطلبات العربية، وإعادة الطلب بالإنجليزية تكفي للتحقق
  \item[\textbf{ب)}] الاستدلال صحيح، لأن الأدوات الحديثة لا تكتب بثقة إلا إذا كانت متأكدة، وجمال الأسلوب يدل على أنها قرأت المصدر الأصلي
  \item[\textbf{ج)}] الخطأ أنه سأل أداة واحدة فقط، ولو سأل أداة ثانية فأعطت الجواب نفسه لصار ذلك دليلًا قاطعًا لا يحتاج إلى أي مصدر آخر
  \item[\textbf{د)}] الخطأ هنا انحياز خوارزمي سببه قلة صور الشعراء في البيانات، والحل الوحيد هو إعادة تدريب الأداة على صور الشعراء القدامى
  \item[\textbf{هـ)}] الأداة قد تكتب معلومة خاطئة بثقة، وهي الهلوسة؛ فالثقة وجمال الأسلوب لا يثبتان الصحة، والدليل هو الرجوع إلى الديوان المحقق
\end{itemize}
\end{minipage}\par\vspace{8pt plus 6pt}
\noindent\begin{minipage}{\linewidth}
\hangindent=0.7cm\hangafter=1\makebox[0.7cm][r]{\textbf{12.}}\textbf{حقق نموذج لتصحيح الإعراب دقة شبه كاملة على أمثلة التدريب، لكنه أخطأ كثيرًا عند تجربته على جمل جديدة من نصوص أدبية قديمة. أي تشخيص وأي علاج هما الأنسب؟}
\begin{itemize}[leftmargin=0.95cm, labelsep=4pt, itemsep=1pt, topsep=3pt, parsep=0pt]
  \item[\textbf{أ)}] السبب أن بيانات الاختبار كانت متنوعة أكثر من اللازم؛ والعلاج أن نختار للاختبار جملًا تشبه أمثلة التدريب قدر الإمكان
  \item[\textbf{ب)}] السبب أن عدد الأمثلة كبير جدًا؛ فنحذف نصف البيانات، لأن النموذج يتعلم أفضل من الأمثلة القليلة مهما كان نوعها وتنوعها
  \item[\textbf{ج)}] حفظ النموذج أمثلة التدريب ولم يكتشف أنماطًا عامة، وبياناته لا تمثّل النصوص القديمة؛ نضيف أمثلة متنوعة ونختبره بأمثلة جديدة
  \item[\textbf{د)}] المشكلة في مرحلة النتيجة فقط؛ نعدّل طريقة عرض النتائج على الشاشة دون المساس بالبيانات أو بالتدريب أو بطريقة الاختبار
  \item[\textbf{هـ)}] النموذج ممتاز والخلل في الجمل الجديدة لأنها صعبة؛ نكتفي باختباره مستقبلًا على أمثلة التدريب نفسها لأنها تعطي النتيجة الحقيقية
\end{itemize}
\end{minipage}\par\vspace{8pt plus 6pt}
\noindent\begin{minipage}{\linewidth}
\hangindent=0.7cm\hangafter=1\makebox[0.7cm][r]{\textbf{13.}}\textbf{حلّل الأمر: «تصرّف كأستاذ جامعي في النقد الأدبي. حلّل قصيدة المساء لخليل مطران لطلبة السنة الثانية ليسانس، في أقل من 300 كلمة، في ثلاثة محاور: الموضوع، الصورة، الإيقاع». أي تحليل صحيح؟}
\begin{itemize}[leftmargin=0.95cm, labelsep=4pt, itemsep=1pt, topsep=3pt, parsep=0pt]
  \item[\textbf{أ)}] الدور: طلبة الثانية ليسانس؛ المهمة: أقل من 300 كلمة؛ السياق: أستاذ في النقد؛ القيود: ثلاثة محاور؛ التنسيق: تحليل القصيدة
  \item[\textbf{ب)}] الدور: أستاذ في النقد؛ المهمة: ثلاثة محاور؛ السياق: أقل من 300 كلمة؛ القيود: طلبة الثانية ليسانس؛ التنسيق: تحليل القصيدة
  \item[\textbf{ج)}] الدور: تحليل القصيدة؛ المهمة: أستاذ في النقد؛ السياق: ثلاثة محاور؛ القيود: طلبة الثانية ليسانس؛ التنسيق: أقل من 300 كلمة
  \item[\textbf{د)}] الدور: أستاذ في النقد؛ المهمة: تحليل القصيدة؛ السياق: طلبة الثانية ليسانس؛ القيود: أقل من 300 كلمة؛ التنسيق: ثلاثة محاور
  \item[\textbf{هـ)}] الدور: أستاذ في النقد؛ المهمة: تحليل القصيدة؛ السياق: ثلاثة محاور؛ القيود: طلبة الثانية ليسانس؛ التنسيق: أقل من 300 كلمة
\end{itemize}
\end{minipage}\par\vspace{8pt plus 6pt}
\noindent\begin{minipage}{\linewidth}
\hangindent=0.7cm\hangafter=1\makebox[0.7cm][r]{\textbf{14.}}\textbf{كتب طالب في آخر بحثه: «استعملتُ أداة ذكاء اصطناعي». لكنه نسخ منها فصلين كاملين دون مراجعة، وفيهما معلومات خاطئة. أي تقييم يطبّق مبادئ الدرس كلها؟}
\begin{itemize}[leftmargin=0.95cm, labelsep=4pt, itemsep=1pt, topsep=3pt, parsep=0pt]
  \item[\textbf{أ)}] عمله سليم تمامًا، لأن الإفصاح يرفع عنه المسؤولية وينقلها إلى الأداة، فالمعلومات الخاطئة أخطاء الأداة وليست أخطاءه هو
  \item[\textbf{ب)}] أفصح لكنه لم يحقق الشفافية لأنه لم يوضّح كيف استعملها، وأخلّ بالمسؤولية لعدم التحقق، ونسخ فصلين باسمه يبقى انتحالًا
  \item[\textbf{ج)}] عمله انتحال فقط، أما المعلومات الخاطئة فلا علاقة لها بالأخلاقيات، لأن التحقق من المعلومات مسألة تقنية لا أخلاقية
  \item[\textbf{د)}] أخلّ بالإفصاح فقط لأن جملته قصيرة جدًا، ولو كتب جملة أطول لأصبح عمله مقبولًا حتى مع نسخ الفصلين ووجود الأخطاء فيهما
  \item[\textbf{هـ)}] لم يخلّ بأي مبدأ، لأن نسخ فصلين نسبة صغيرة، والدرس لا يعدّ انتحالًا إلا نسخ البحث كله من أوله إلى آخره دون استثناء
\end{itemize}
\end{minipage}\par\vspace{8pt plus 6pt}
\noindent\begin{minipage}{\linewidth}
\hangindent=0.7cm\hangafter=1\makebox[0.7cm][r]{\textbf{15.}}\textbf{وجد طالب ثلاثة مصادر حول أثر الألعاب التعليمية على التحصيل: (1) مقال محكّم حديث بمراجع، (2) صفحة شركة تبيع لعبة تعليمية وتَعِد بنتائج مضاعفة، (3) ملخص أنتجته أداة ذكاء اصطناعي دون مراجع. أي تقييم صحيح؟}
\begin{itemize}[leftmargin=0.95cm, labelsep=4pt, itemsep=1pt, topsep=3pt, parsep=0pt]
  \item[\textbf{أ)}] يرفض المصادر الثلاثة جميعًا، لأن الموضوع حديث، ولا يمكن الوثوق بأي مصدر لم يمر على نشره أكثر من عشرين سنة كاملة
  \item[\textbf{ب)}] يعتمد المقال المحكّم؛ ويحذر من صفحة الشركة لأن هدفها البيع ولا مراجع لها؛ ويستعمل الملخص للفهم فقط لأنه ليس مصدرًا علميًا
  \item[\textbf{ج)}] يعتمد صفحة الشركة لأنها الأكثر خبرة بمنتجها، ويعدّ المقال المحكّم والملخص مصدرين متساويين في درجة الثقة والقيمة العلمية
  \item[\textbf{د)}] يعتمد ملخص الأداة أولًا لأنه الأحدث والأوضح، ثم صفحة الشركة لأنها متخصصة في الألعاب، ويترك المقال لأنه صعب ومعقد
  \item[\textbf{هـ)}] يعتمد المصادر الثلاثة بالدرجة نفسها، لأن معايير التقييم تخص الكتب الورقية فقط، ولا تنطبق على المقالات والمواقع والأدوات
\end{itemize}
\end{minipage}\par\vspace{8pt plus 6pt}
\noindent\begin{minipage}{\linewidth}
\hangindent=0.7cm\hangafter=1\makebox[0.7cm][r]{\textbf{16.}}\textbf{طلب أستاذ من أداة متعددة الأنماط تحويل درس «الجملة الاسمية» إلى مخطط مفاهيمي، فجاء المخطط: الجملة الاسمية $\leftarrow$ المبتدأ (اسم منصوب)، الخبر (مفرد فقط)، الفاعل. ما الأخطاء التي تكشفها أسئلة المراجعة؟}
\begin{itemize}[leftmargin=0.95cm, labelsep=4pt, itemsep=1pt, topsep=3pt, parsep=0pt]
  \item[\textbf{أ)}] خطأ واحد فقط هو إضافة الفاعل، أما المبتدأ المنصوب والخبر المفرد فهما صحيحان تمامًا حسب المخطط الذي عرضه الدرس
  \item[\textbf{ب)}] خطأ واحد فقط في الشكل: الأسهم غير ملوّنة، أما المفاهيم والعلاقات فصحيحة، لأن مراجعة المخطط تقتصر على الألوان والخطوط
  \item[\textbf{ج)}] خطأ واحد فقط هو الخبر المفرد، لأن الخبر لا يكون مفردًا أبدًا، أما المبتدأ المنصوب والفاعل فهما من مفاهيم الجملة الاسمية
  \item[\textbf{د)}] لا أخطاء في المخطط، لأن الأداة متعددة الأنماط لا تخطئ في تحويل النص إلى مخطط، وإنما تخطئ فقط في قراءة الصور والأرقام
  \item[\textbf{هـ)}] علاقة خاطئة: المبتدأ مرفوع لا منصوب؛ ونقص فكرة أساسية: الخبر قد يكون جملة أو شبه جملة؛ ومفهوم مضاف غير موجود: الفاعل
\end{itemize}
\end{minipage}\par\vspace{8pt plus 6pt}
\noindent\begin{minipage}{\linewidth}
\hangindent=0.7cm\hangafter=1\makebox[0.7cm][r]{\textbf{17.}}\textbf{طلب طالب من أداة تلخيص مقال من عشر صفحات، فجاء الملخص منظمًا، لكنه أغفل النتيجة الرئيسية للمقال وأضاف إحصائية غير موجودة فيه. أي معالجة تطبّق خطوات التلخيص ومراجعته؟}
\begin{itemize}[leftmargin=0.95cm, labelsep=4pt, itemsep=1pt, topsep=3pt, parsep=0pt]
  \item[\textbf{أ)}] يقدّم الملخص كما هو دون ذكر المصدر، لأن الأداة كتبته بأسلوبها، فلا يحتاج إلى ذكر المقال الأصلي في قائمة المراجع
  \item[\textbf{ب)}] يقارن الملخص بالنص الأصلي فيكتشف الحذف والإضافة، ثم يعيد التلخيص: يحدد الأفكار الأساسية، ويكتب بأسلوبه، ويذكر المصدر
  \item[\textbf{ج)}] يحذف الإحصائية المضافة فقط ويعتمد الباقي، لأن حذف النتيجة الرئيسية لا يؤثر في المعنى ما دامت بقية الأفكار موجودة
  \item[\textbf{د)}] يعتمد الملخص لأنه منظم وواضح، فالتنظيم الجيد دليل على أن الأداة فهمت المقال، والإحصائية المضافة تزيد البحث قيمة
  \item[\textbf{هـ)}] يطلب من الأداة ملخصًا أطول ويعتمده دون قراءة المقال، لأن الملخص الأطول يحتوي بالضرورة على كل الأفكار الأساسية
\end{itemize}
\end{minipage}\par\vspace{8pt plus 6pt}
\noindent\begin{minipage}{\linewidth}
\hangindent=0.7cm\hangafter=1\makebox[0.7cm][r]{\textbf{18.}}\textbf{صمّمت مؤسسة أداة لاقتراح نصوص قراءة لتلاميذ كل الولايات، ودُرّبت على نصوص من بيئة الشمال فقط. اشتكى أساتذة الجنوب من أمثلة غريبة عن تلاميذهم، فردّ المصمم: «الأداة هي التي اختارت النصوص، لا أنا». أي تحليل يجمع مفاهيم الفصلين الثاني والرابع بدقة؟}
\begin{itemize}[leftmargin=0.95cm, labelsep=4pt, itemsep=1pt, topsep=3pt, parsep=0pt]
  \item[\textbf{أ)}] انحياز خوارزمي سببه بيانات غير ممثِّلة؛ والمسؤولية على الإنسان لا على الأداة، فعلى المصمم تنويع البيانات واختبارها في كل البيئات
  \item[\textbf{ب)}] هلوسة سببها صياغة الطلب؛ والمسؤولية على الأساتذة لأنهم لم يكتبوا أوامر دقيقة، فعليهم تحسين أوامرهم بإضافة الدور والسياق
  \item[\textbf{ج)}] انتحال لأن الأداة نسخت نصوص الشمال؛ والحل أن تصرّح المؤسسة بذلك في جملة إفصاح، فيصبح الاستعمال مقبولًا دون أي تعديل
  \item[\textbf{د)}] لا يوجد انحياز لأن عدد النصوص كبير؛ والمسؤولية على الأداة كما قال المصمم، لأنها هي التي اختارت النصوص وليس الإنسان
  \item[\textbf{هـ)}] انحياز خوارزمي لا أحد مسؤول عنه، لأن الانحياز صفة طبيعية في كل نموذج، ولا يمكن التقليل منه بأي طريقة من الطرق
\end{itemize}
\end{minipage}\par\vspace{8pt plus 6pt}
\noindent\begin{minipage}{\linewidth}
\hangindent=0.7cm\hangafter=1\makebox[0.7cm][r]{\textbf{19.}}\textbf{أراد أستاذ تحويل صورة صفحة مخطوطة إلى نص رقمي، ثم إلى مخطط مفاهيمي لتلاميذ السنة الثالثة متوسط. أي خطة عمل تجمع بين هندسة الأوامر وخطوات العمل مع الأنماط بأفضل شكل؟}
\begin{itemize}[leftmargin=0.95cm, labelsep=4pt, itemsep=1pt, topsep=3pt, parsep=0pt]
  \item[\textbf{أ)}] تحويل الصورة إلى صوت أولًا لتسهيل القراءة، ثم كتابة أمر طويل جدًا بكل التفاصيل في خطوة واحدة، ثم قبول أول إجابة تظهر
  \item[\textbf{ب)}] طلب المخطط مباشرة من الصورة دون المرور بالنص، ثم مقارنة المخطط بالصورة من حيث الألوان فقط، دون مراجعة المفاهيم والعلاقات
  \item[\textbf{ج)}] كتابة أمر واحد يجمع النقل والمخطط والترجمة والتلخيص معًا، لأن جمع المهام في أمر واحد يوفّر الوقت ويحسّن جودة الإجابة
  \item[\textbf{د)}] صورة واضحة مع أمر دقيق لنقلها نصًا، ثم مقارنة النص بالمخطوطة، ثم أمر يحدد المهمة والسياق والتنسيق للمخطط، ثم مراجعة المخطط
  \item[\textbf{هـ)}] رفع أي صورة متاحة مع كلمة «حوّل»، ثم اعتماد النص الناتج مباشرة، ثم طلب «مخطط» دون تحديد، ونشر المخطط دون أي مراجعة
\end{itemize}
\end{minipage}\par\vspace{8pt plus 6pt}
\noindent\begin{minipage}{\linewidth}
\hangindent=0.7cm\hangafter=1\makebox[0.7cm][r]{\textbf{20.}}\textbf{طلب طالب من أداة «قائمة مراجع حول تعليم البلاغة»، فقدّمت عشرة مراجع منسقة، ثلاثة منها غير موجودة أصلًا. نقلها كلها إلى بحثه دون تحقق ودون إفصاح. أي تحليل يجمع مفاهيم الفصول الأول والرابع والخامس بدقة؟}
\begin{itemize}[leftmargin=0.95cm, labelsep=4pt, itemsep=1pt, topsep=3pt, parsep=0pt]
  \item[\textbf{أ)}] المراجع الوهمية هلوسة؛ ومحتوى الأداة ليس مصدرًا علميًا فيجب التحقق من كل مرجع؛ وعدم التحقق والإفصاح إخلال بالمسؤولية والإفصاح
  \item[\textbf{ب)}] المشكلة الوحيدة غياب الإفصاح؛ فلو كتب جملة إفصاح لأصبحت المراجع الوهمية مقبولة، لأن الإفصاح ينقل المسؤولية إلى الأداة
  \item[\textbf{ج)}] المشكلة الوحيدة أن العدد عشرة كبير؛ فلو طلب ثلاثة مراجع فقط لكانت كلها صحيحة، لأن الأداة لا تهلوس مع القوائم القصيرة
  \item[\textbf{د)}] لا مشكلة في المراجع لأنها منسقة، والتنسيق الجيد دليل على وجودها؛ والطالب غير مسؤول لأن الأداة هي التي اقترحت القائمة
  \item[\textbf{هـ)}] المراجع الوهمية انحياز خوارزمي؛ ومحتوى الأداة أعلى المصادر ثقة فلا حاجة إلى التحقق؛ والإفصاح غير مطلوب في قائمة المراجع
\end{itemize}
\end{minipage}\par\vspace{8pt plus 6pt}
\end{multicols}
\vfill
\begin{center}
{\khatt\textbf{جدول الإجابات}}\\[3pt]
\renewcommand{\arraystretch}{1.6}\setlength{\tabcolsep}{2.5pt}
\begin{tabular}{|>{\centering\arraybackslash}m{1.9cm}|*{20}{>{\centering\arraybackslash}m{0.5cm}|}}
\hline
\textbf{رقم السؤال} & \textbf{1} & \textbf{2} & \textbf{3} & \textbf{4} & \textbf{5} & \textbf{6} & \textbf{7} & \textbf{8} & \textbf{9} & \textbf{10} & \textbf{11} & \textbf{12} & \textbf{13} & \textbf{14} & \textbf{15} & \textbf{16} & \textbf{17} & \textbf{18} & \textbf{19} & \textbf{20} \\ \hline
\textbf{الإجابة} &  &  &  &  &  &  &  &  &  &  &  &  &  &  &  &  &  &  &  &  \\[8pt] \hline
\end{tabular}

\vspace{8pt}
\renewcommand{\arraystretch}{2.2}
\begin{tabular}{|>{\centering\arraybackslash}m{7.5cm}|>{\centering\arraybackslash}m{3.2cm}|>{\centering\arraybackslash}m{3.2cm}|}
\hline
\textbf{الاسم واللقب} & \textbf{الفوج} & \textbf{العلامة} \\ \hline
 & & \hfill /20 \\ \hline
\end{tabular}
\end{center}
\end{document}
